\section{相对观测的可辨识性}
\label{sec:smap-identifiability}

在计算任何坐标或映射之前，先要弄清观测能够确定什么。一组测量即使完全没有误差，也可能对某些量毫无约束：把整幅图平移一下，所有距离都不变，距离表当然无法告诉我们原点在哪里。本节先区分两类相对观测，再把观测不能辨认的自由度分成两种：一种来自坐标系的任意选择，另一种来自观测本身的缺失。

\paragraph{两类相对观测}

第一类观测只涉及点与点之间的关系。给定点标签集合$V=\{1,\ldots,n\}$和已观测的点对集合$E$，对每个$\{i,j\}\in E$测得一个距离$d_{i,j}\geq0$。未知的是坐标$\bm x_1,\ldots,\bm x_n\in\mathbb R^d$，要求它们满足
\[
 \lVert\bm x_i-\bm x_j\rVert_2=d_{i,j},\qquad\{i,j\}\in E .
\]
$E$包含全部点对时称距离表完整，否则称为部分距离。这里的数据只有标签与数，没有任何外部坐标。

第二类观测涉及同一组物理点在两种坐标表示下的记录。对每个$i$，已知源坐标$\bm x_i\in\mathbb R^d$与目标坐标$\bm y_i\in\mathbb R^m$，并且知道它们描述同一个点。未知的是把源坐标送到目标坐标的映射$f$，要求$f(\bm x_i)\approx\bm y_i$。两类问题的未知量不同：前者的未知量是$n$个点的坐标，共$nd$个数；后者的未知量是一个映射，点越多，对它的约束越多。

\paragraph{整体欧氏自由度}

先看距离不能辨认什么。取正交矩阵$\bm Q\in\mathrm O(d)$和向量$\bm b\in\mathbb R^d$，令$\bm x_i'=\bm Q\bm x_i+\bm b$。差向量变为$\bm x_i'-\bm x_j'=\bm Q(\bm x_i-\bm x_j)$，正交矩阵保持长度，所以每一对点的距离都不变。也就是说，平移、旋转和反射作用于整组点以后，任何距离观测都看不出区别。

\begin{definition}[合同的点配置]
\label{def:smap-congruent}
两组点$\bm x_1,\ldots,\bm x_n$与$\bm x_1',\ldots,\bm x_n'\in\mathbb R^d$称为\term{合同}（congruent），如果存在$\bm Q\in\mathrm O(d)$和$\bm b\in\mathbb R^d$，使$\bm x_i'=\bm Q\bm x_i+\bm b$对所有$i$成立。
\end{definition}

由上一段的计算，合同的配置给出同一张距离表。因此距离观测最多只能把点配置确定到合同的程度，这一限制与测量是否完整、是否精确都无关。平移有$d$个参数，$\mathrm O(d)$的每个连通分支有$d(d-1)/2$个参数，所以合同带来$d(d+1)/2$个连续自由度，外加一个反射的二选一。平面上是$3$个连续自由度，空间中是$6$个。

反过来，完整距离表是否把配置确定到合同为止？第~\ref{sec:smap-full-distance}节会证明答案是肯定的：两组点的全部距离相同，它们就合同。这说明对完整距离表而言，合同是唯一的不确定性，坐标重构的任务是在合同意义下找出配置。

\paragraph{缺失观测的形变自由度}

观测不完整时，不确定性可能多于合同。在平面上取单位正方形的四个顶点，只测四条边长，都等于$1$，不测两条对角线。把正方形压成内角为$\theta$的菱形，四条边长仍是$1$，但两条对角线长分别变为$2\sin(\theta/2)$与$2\cos(\theta/2)$。$\theta\neq\pi/2$时对角线与正方形不同，所以菱形与正方形不合同，却满足全部观测。

这里的自由度来自观测不足，可以用计数来理解。平面上$n$个点有$2n$个坐标，扣去合同的$3$个自由度，形状本身还有$2n-3$个自由度；$n=4$时为$5$个，四个边长约束一般只能消去其中四个，剩下的一个就是菱形的内角。计数只给出必要条件，约束够多也未必确定形状，第~\ref{sec:smap-partial-distance}节将用具体例子说明这一点。

\paragraph{固定参照的作用}

为了得到一组具体坐标，常常人为固定参照。在平面上可以规定$\bm x_1=\bm0$，$\bm x_2$位于第一坐标轴正半轴上，$\bm x_3$位于上半平面。这三条规定恰好消去平移、旋转和反射：任何配置都与唯一一个满足规定的配置合同，只要$\bm x_1,\bm x_2,\bm x_3$不共线。

但规定只能消去约定上的自由度。回到菱形的例子，取$\bm x_1=\bm0$、$\bm x_2=(1,0)^\top$，并把第四个点放在上半平面，内角$\theta$仍然可以连续变化，四个边长仍然成立。形变自由度是由缺失的观测造成的，再多的坐标约定也补不出缺失的对角线长度。

还要区分约定与观测。若某些点的绝对位置是用仪器测出来的，例如卫星定位给出的锚点坐标，它们就是额外的观测，可能与距离数据相互矛盾，需要进入拟合残差；只有任意选取、不影响距离的规定，才是消去合同自由度的约定。第二类问题中也有对应的区分：两套坐标系都已给定，映射本身没有约定上的自由度，但必须事先选定映射族，第~\ref{sec:smap-correspondence-models}节将说明这一点。
